\section{Differentiabilitet}

\subsection{Differentiabilitet af skalare funktioner}

En funktion \emph{fra en åben mængde} er differentiabel i $x_0$ hvis der eksisterer
en Epsilon-funktion $\varepsilon(h)$ sådan at ligningen går op:

\[
    f(x_0+h) - f(x_0) = f'(x_0) \cdot h + |h| \cdot \varepsilon(h), \qquad \forall h \in \mathbb{R}
    \kildetag{Lemma 3.6.1}
\]


Skalare funktioner af flere variable:
\[
    f(\vec x_0 + \vec h) - f(\vec x_0) = (\nabla f(\vec x_0))^T \cdot \vec h + ||\vec h|| \cdot \varepsilon(\vec h), \qquad \forall \vec h \in \mathbb{R}^n
    \kildetag{Definition 3.6.1}
\]

\subsubsection{Epsilon-funktioner}
En Epsilon-funktion $\varepsilon(h)$ er en vilkårlig funktion der går mod 0 når $h$ går mod 0.

\eksempel{
    Vis at $f: \mathbb R^2 \to \mathbb R, f(x,y) = x^2 + y^2$ er differentiabel i hele $\mathbb R^2$.

    Hvis vi kan vise at denne ligning går op er det vist:
    \[
        f(\vec x_0 + \vec h) - f(\vec x_0) = (\nabla f(x_0,y_0))^T \cdot \vec h + ||\vec h|| \cdot \varepsilon(\vec h)
    \]
    Vi kan udvide vektorerne og finde gradienten:
    \[
        f(x_0 + h_x, y_0 + h_y) - f(x_0, y_0) = (\nabla f(x_0,y_0))^T \cdot 
        \begin{bmatrix}
            h_x &
            h_y
        \end{bmatrix}^T + ||\vec h|| \cdot \varepsilon(\vec h)
    \]
    \[
        \nabla f(x_0,y_0) = (2x_0, 2y_0)
    \]
    Udregner funktionsværdierne, indsætter gradienten og regner $||\vec h||$:
    \[
        (x_0 + h_1)^2 + (y_0 + h_2)^2 - (x_0^2 + y_0^2) = 
        \begin{bmatrix}
            2x_0 &
            2y_0
        \end{bmatrix}        
        \cdot 
        \begin{bmatrix}
            h_1 &
            h_2
        \end{bmatrix}^T 
        + \sqrt{h_1^2 + h_2^2}
        \cdot \varepsilon(\vec h)
    \]
    Udvider de første 2 led, og ganger de 2 vektorer sammen:
    \[
       ( x_0^2 + 2x_0h_1 + h_1^2) + (y_0^2 + 2y_0h_2 + h_2^2) - (x_0^2 + y_0^2) = 2x_0h_1 + 2y_0h_2 + \sqrt{h_1^2 + h_2^2} \cdot \varepsilon(\vec h)
    \]
    Reducerer udtrykket:
    \[
        h_1^2 + h_2^2 = \sqrt{h_1^2 + h_2^2} \cdot \varepsilon(\vec h)
    \]
    Deler med $\sqrt{h_1^2 + h_2^2}$
    \[
        \frac{h_1^2+h_2^2}{\sqrt{h_1^2 + h_2^2}} = \sqrt{h_1^2 + h_2^2} = \varepsilon(\vec h)
    \]

    Det kan nu ses at $\vec h \to \vec 0 \implies \varepsilon(\vec h) \to \vec 0$, dermed bevist.
}

\subsubsection{Sammensatte funktioner differentieret}

\begin{align*}
(f+g)'(x) &= f'(x) + g'(x) \\
(f \cdot g)'(x) &= f'(x) \cdot g(x) + f(x) \cdot g'(x) \\
(f \circ g)'(x) &= f'(g(x)) \cdot g'(x)
\end{align*}

\subsubsection{Tangent i punkt}
Tangentvektoren til en vektorfunktion $\vec r$ i punktet $\vec t_0$ findes ved
\[
r(t_0) + s \cdot r'(t_0), \quad s \in \mathbb{R}
\]

\subsubsection{Differentiabilitet af skalare funktioner af flere variable}
Hvis $f: U \to R$ er differentiabel i $x_0$, så findes alle partielt afledte i $x_0$.
\kilde{Theorem 3.6.2}

Hvis alle partielt afledte er kontinuerte i $x_0$, så er $f$ differentiabel i $x_0$.
\kilde{Theorem 3.6.3}

\subsection{Partielt afledt}
Antag alle andre variable er konstanter og differentier som normalt.

\eksempel{
    Den partielt afledte af $f(x,y)= x^2y + 3xy^2$ med hensyn til $x$ er
    \[
        \frac{\partial f}{\partial x} = 2xy + 3y^2
    \]
    Med hensyn til $y$ er den partielt afledte
    \[
        \frac{\partial f}{\partial y} = x^2 + 6xy
    \]
}

\subsubsection{Dobbelt-afledt}

\[
    \frac{\partial^2 f}{\partial x \partial y} = \frac{\partial}{\partial x} \left( \frac{\partial f}{\partial y} \right)
\]

\eksempel{
    Find den dobbelt-afledte af $f(x,y) = \sin(x^2 + y)$ mhht $x$.
    \begin{align*}
        \frac{\partial^2 f}{\partial x^2} \left( \sin(x^2 + y) \right) &= \frac{\partial f}{\partial x} \left( \frac{\partial f}{\partial x} \left( \sin(x^2+y) \right) \right) \\
        &= \frac{\partial f}{\partial x} \left( \cos(x^2 + y) \cdot 2x \right)
    \end{align*}

    \begin{align*}
        \frac{\partial f}{\partial x} (\cos(x^2 + y) \cdot 2x)
        &= (\cos(x^2 + y))' \cdot 2x &+& \cos(x^2 + y) \cdot (2x)' \\
        &= -2x \sin(x^2 + y) \cdot 2x &+& \cos(x^2 + y) \cdot 2 \\
        &= -4x^2 \sin(x^2 + y) &+& 2 \cos(x^2 + y)
    \end{align*}
}

\subsubsection{Hesse for kvadratisk form}
Hesse matricen for en kvadratisk form $f(\vec x) = \vec x^T \mtrix A \vec x + \vec x^T \vec b + c$ er altid $H_f = 2 \mtrix A$.

\subsection{Gradientvektor}
\[
\nabla f (x,y) = \left( \frac{\partial f}{\partial x}, \frac{\partial f}{\partial y} \right)
\]

Står vinkelret på niveakurverne til $f$ og har retning af største stigning i $f$.

\eksempel{
    Gradientvektoren til $f(x,y) = x^2y + 3xy^2$ er
    \[
        \nabla f(x,y) = (2xy + 3y^2, x^2 + 6xy)
    \]
}

\subsection{Retnings-afledt}
Den retnings-afledte af $f$ i retningen af en vektor $\vec v$ er givet ved
\[
\nabla_{\vec v} f(x,y) = \inner{\nabla f(x,y)}{\frac{\vec v}{||\vec v||}}
\]

Den retnings-afledte måler den hastighed $f$ ændrer sig i retningen af $\vec v$.

\subsection{Gradient Ascent og Gradient Descent}
Find lokal maksimum: Tag startpunkt, bevæg dig i retning af gradienten, gentag.

Find lokal minimum: Tag startpunkt, bevæg dig i den modsatte retning af gradienten, gentag.

\subsection{Hesse Matricer}
\label{sec:HesseMatricer}
Hesse matricen af en funktion $f(\vec x)$ er en kvadratisk matrix, der indeholder de andenordens partielt afledte:
\begingroup
\renewcommand*{\arraystretch}{1.5}
\[
    H_f(x,y) = \begin{bmatrix}
        \frac{\partial^2 f}{\partial x^2} & \frac{\partial^2 f}{\partial x \partial y} \\
        \frac{\partial^2 f}{\partial y \partial x} & \frac{\partial^2 f}{\partial y^2}
    \end{bmatrix}
    \qquad
    H_f(x,y,\ldots,z) = \begin{bmatrix}
        \frac{\partial^2 f}{\partial x^2} & \frac{\partial^2 f}{\partial x \partial y} & \cdots & \frac{\partial^2 f}{\partial x \partial z} \\
        \frac{\partial^2 f}{\partial y \partial x} & \frac{\partial^2 f}{\partial y^2} & \cdots & \frac{\partial^2 f}{\partial y \partial z} \\
        \vdots & \vdots & \ddots & \vdots \\
        \frac{\partial^2 f}{\partial z \partial x} & \frac{\partial^2 f}{\partial z \partial y} & \cdots & \frac{\partial^2 f}{\partial z^2}
    \end{bmatrix}
\]
\endgroup

\bemærkBig{
    \( \frac{\partial^2 f}{\partial y \partial x} = \frac{\partial^2 f}{\partial x \partial y} \) for kontinuerete funktioner.
    \kilde{Clairaut's sætning}

    $\implies$ Hesse matricen af en kontinuert funktion er symmetrisk.
}










\subsection{Jacobi Matricer}
Givet vektorfunktion $\vec f: \mathbb{R}^n \to \mathbb{R}^m$, er Jacobi matricen $J_{\vec f}$ en $m \times n$ matrix, der indeholder de førsteordens partielt afledte:
\begingroup
\renewcommand*{\arraystretch}{1.5}
\[
    J_{\vec f}(\vec x) = \begin{bmatrix}
        \frac{\partial f_1}{\partial x_1} & \frac{\partial f_1}{\partial x_2} & \cdots & \frac{\partial f_1}{\partial x_n} \\
        \frac{\partial f_2}{\partial x_1} & \frac{\partial f_2}{\partial x_2} & \cdots & \frac{\partial f_2}{\partial x_n} \\
        \vdots & \vdots & \ddots & \vdots \\
        \frac{\partial f_m}{\partial x_1} & \frac{\partial f_m}{\partial x_2} & \cdots & \frac{\partial f_m}{\partial x_n}
    \end{bmatrix}
\]
\endgroup

\subsubsection{Jacobianten}
Jacobianten er determinanten af Jacobi matricen.

\subsection{Kædereglen}

\subsubsection{Kædereglen for skalare funktioner af flere variable}

Givet $f: \mathbb{R}^n \to \mathbb{R}$ og $\vec g: \mathbb{R}^m \to \mathbb{R}^n$:
\[
    (f \circ \vec g)'(\vec x) = \nabla f(\vec g(\vec x))^T \cdot \vec g'(\vec x)
\]

\kilde{Korollar 3.7.2}

\eksempel{
    Givet $f(x,y) = x^2 + y^2, \vec g(y) = (u \cos(u), u \sin(u))$, find $(f \circ \vec g)'(u)$.

    \[
        \nabla f(x,y) = 
        \begin{bmatrix}
            2x & 2y
        \end{bmatrix}^T
        \implies
        \nabla f(\vec g(u)) = 
        \begin{bmatrix}
            2u \cos(u) & 2u \sin(u)
        \end{bmatrix}^T
    \]

    \[
        \vec g'(u) = 
        \begin{bmatrix}
            \cos(u) - u \sin(u) \\
            \sin(u) + u \cos(u)
        \end{bmatrix}
    \]

    \begin{align*}        
        &\begin{bmatrix}
            2u \cos(u) & 2u \sin(u)
        \end{bmatrix}^T
        \cdot
        \begin{bmatrix}
            \cos(u) - u \sin(u) \\
            \sin(u) + u \cos(u)
        \end{bmatrix} \\
        &= 2u \cos^2(u) - 2u^2 \cos(u) \sin(u) + 2u \sin^2(u) + 2u^2 \sin(u) \cos(u) \\
        &= 2u (\cos^2(u) + \sin^2(u)) \\
        &= 2u
    \end{align*}

    
}

\subsubsection{Kædereglen med Jacobi}
Givet $\vec f: \mathbb{R}^n \to \mathbb{R}^m$ og $\vec g: \mathbb{R}^m \to \mathbb{R}^l$, så er kædereglen for sammensætningen $\vec h = \vec g \circ \vec f$ givet ved:
\[
    J_{\vec h}(\vec x) = J_{\vec g}(\vec f(\vec x)) \cdot J_{\vec f}(\vec x)
\]

\kilde{Theorem 3.8.4}

\eksempel{Givet funktioner $f$ og $g$, find Jacobi matricen for $h = f \circ g$.
\[
f: \mathbb R^2 \to \mathbb R^2, f(x,y) = 
\begin{bmatrix}
x^2 & x+y
\end{bmatrix}^T
\qquad
g: \mathbb R^2 \to \mathbb R^2, g(r,\theta) = \begin{bmatrix}
r \cos (\theta) & r \sin (\theta)
\end{bmatrix}^T
\]

\textbf{1.} Find $J_f(x,y)$ og $J_g(r,\theta)$
\begin{align*}
J_f(x,y) = \begin{bmatrix}
\frac{\partial f_1}{\partial x} & \frac{\partial f_1}{\partial y} \\
\frac{\partial f_2}{\partial x} & \frac{\partial f_2}{\partial y}
\end{bmatrix}
&=
\begin{bmatrix}
2x & 0 \\
1 & 1
\end{bmatrix}
\\
J_g(r,\theta) = \begin{bmatrix}
\frac{\partial g_1}{\partial r} & \frac{\partial g_1}{\partial \theta} \\
\frac{\partial g_2}{\partial r} & \frac{\partial g_2}{\partial \theta}
\end{bmatrix}
&=
\begin{bmatrix}
\cos(\theta) & -r \sin(\theta) \\
\sin(\theta) & r \cos(\theta)
\end{bmatrix}
\end{align*}

\textbf{2.} Find $J_f(g(r,\theta))$
Her skal vi tage $(x,y)$ som input til $g$, hvor vi får $g(x,y) = (r \cos(\theta), r \sin(\theta))$. Vi udskifter så $x$ og $y$ i $J_f$ med $r \cos(\theta)$ og $r \sin(\theta)$:
\[
J_f(g(r,\theta)) = 
\begin{bmatrix}
2r \cos(\theta) & 0 \\
1 & 1
\end{bmatrix}
\]

\textbf{3.} Multiplicer matricerne

\[
J_h =
\begin{bmatrix}
2r \cos(\theta) & 0 \\
1 & 1
\end{bmatrix}
\cdot
\begin{bmatrix}
\cos(\theta) & -r \sin(\theta) \\
\sin(\theta) & r \cos(\theta)
\end{bmatrix}
=
\begin{bmatrix}
2r \cos^2(\theta) & -2r^2 \cos(\theta) \sin(\theta) \\
\cos(\theta) + \sin(\theta) & r (\cos(\theta) - \sin(\theta))
\end{bmatrix}
\]

}